\documentclass[12pt,a4paper]{article}

% ==================== 基础包 ====================
\usepackage[UTF8]{ctex}
\usepackage{geometry}
\geometry{top=2.5cm,bottom=2.5cm,left=2.8cm,right=2.8cm}
\usepackage{graphicx}
\usepackage{booktabs}
\usepackage{longtable}
\usepackage{array}
\usepackage{multirow}
\usepackage{enumitem}
\usepackage{xcolor}
\usepackage{titlesec}
\usepackage{fancyhdr}
\usepackage{lastpage}
\usepackage{hyperref}
\usepackage{float}
\usepackage{caption}
\usepackage{subcaption}
\usepackage{colortbl}

% ==================== 页面设置 ====================
\pagestyle{fancy}
\fancyhf{}
\fancyhead[C]{\small 北京印刷学院·"千人百村"实习实践活动成果报告}
\fancyfoot[C]{\thepage}
\renewcommand{\headrulewidth}{0.4pt}

% ==================== 标题格式 ====================
\titleformat{\section}{\Large\bfseries\centering}{\chinese{section}、}{0em}{}
\titleformat{\subsection}{\large\bfseries}{（\chinese{subsection}）}{0.5em}{}
\titleformat{\subsubsection}{\normalsize\bfseries}{\arabic{subsubsection}.}{0.5em}{}

% ==================== 颜色定义 ====================
\definecolor{titleblue}{RGB}{0,82,147}
\definecolor{accentgreen}{RGB}{34,139,34}
\definecolor{lightyellow}{RGB}{255,250,205}
\definecolor{lightgray}{RGB}{245,245,245}

% ==================== 文档开始 ====================
\begin{document}

% ==================== 封面 ====================
\begin{titlepage}
\begin{center}
\vspace*{2cm}
{\small 北京印刷学院·"千人百村"实习实践活动成果报告}

\vspace{3cm}

{\Huge\bfseries\color{titleblue} "千人百村"大学生实习实践活动}\\[0.8cm]
{\Huge\bfseries\color{titleblue} 成\quad 果\quad 报\quad 告}

\vspace{3cm}

\begin{tabular}{rl}
\textbf{学\qquad 校：} & 北京印刷学院 \\[0.5em]
\textbf{专\qquad 业：} & 智能科学与技术 \\[0.5em]
\textbf{实践小队：} & "数智渔韵"实践队 \\[0.5em]
\textbf{实践地点：} & 北京市平谷区马昌营镇 \\[0.5em]
\end{tabular}

\vspace{2cm}

\begin{tabular}{rl}
\textbf{指导教师：} & 蒋祎 \\[0.5em]
\textbf{课题组成员：} & 石宇含\quad 范潇麟\quad 高靖鲲 \\[0.5em]
\end{tabular}

\vspace{2cm}

{\large 2026年7月}

\end{center}
\end{titlepage}

% ==================== 目录 ====================
\newpage
\tableofcontents
\newpage

% ==================== 一、活动背景 ====================
\section{活动背景}

"民族要复兴，乡村必振兴。"在全面推进乡村振兴战略的进程中，首都高校"千人百村"暑期社会实践活动号召广大青年学子深入京郊大地，以专业所学服务乡村发展。北京印刷学院积极响应号召，智能科学与技术专业石宇含、范潇麟、高靖鲲三位同学在指导教师蒋祎的带领下，组建"数智渔韵"实践队，于2026年7月14日至19日赴北京市平谷区马昌营镇开展实地调研与实践活动。

平谷区作为首都生态涵养区，渔业养殖面积达2000余亩，是北京地区重要的淡水鱼供应基地。马昌营镇下辖多个以渔业为主要产业的传统渔村，同时也在积极发展宠物文化产业等新兴产业。本次实践聚焦传统渔业现状调研，同时了解镇域内特色产业生态，探索数智化赋能传统产业转型的可能路径。

% ==================== 二、团队简介 ====================
\section{团队简介}

"数智渔韵"实践队由北京印刷学院智能科学与技术专业的3名大二本科生和1名指导教师组成。团队名称寓意"以数智技术赋能传统渔业"，体现了将人工智能、数据分析等专业知识应用于乡村振兴实践的理念。

指导教师蒋祎负责整体方案设计与学术指导。石宇含、范潇麟、高靖鲲三位同学分别承担实地调研、影像拍摄与记录、数据整理与报告撰写等工作。团队成员具备智能科学与技术专业背景，对人工智能、数据分析、数字化工具应用有较为系统的学习基础，能够在调研过程中以技术视角审视传统产业痛点，提出数智化升级的建议方案。

\begin{figure}[H]
\centering
\fbox{\parbox{0.8\textwidth}{\centering\vspace{2cm}【图1 实践队在天井村调研合影】\vspace{2cm}}}
\caption{实践队在天井村调研合影}
\end{figure}

\begin{figure}[H]
\centering
\fbox{\parbox{0.8\textwidth}{\centering\vspace{2cm}【图2 实践队在马昌营村调研合影】\vspace{2cm}}}
\caption{实践队在马昌营村调研合影}
\end{figure}

% ==================== 三、渔村实地调研 ====================
\section{渔村实地调研}

实践队于2026年7月14日至19日，先后赴北定福村、东双营村、天井村、马昌营村、王各庄村、后芮营村、西海子村七个村庄开展渔业调研。以下按村庄分节记录各村的调研情况。

% ---------- 3.1 北定福村 ----------
\subsection{北定福村渔业调研}

北定福村是实践队调研的第一站。该村鱼塘呈现"集中连片、各户承包"的特征，多个鱼塘紧密相连，由不同养殖户分别承包经营。

\subsubsection{养殖模式与规模}

该村鱼塘为集中连片布局，各户分别承包经营。相比散养村，这一模式已具备了"抱团降成本"的雏形，在组织形式上有一定基础。自己包的，没有合伙的，连片单独（只有一个池塘，很多户包），能省不少成本。

\subsubsection{经营现状与痛点}

养殖户反映的核心问题包括：
\begin{itemize}[leftmargin=2em]
    \item \textbf{鱼价不稳定}：近年来鱼价波动较大，养殖户收入难以保障；
    \item \textbf{成本持续增加}：饲料、电费、药品和人工成本均大幅上涨；
    \item \textbf{利润空间压缩}：鱼价与十年前基本持平，但成本逐年攀升，利润空间被不断压缩。
\end{itemize}

\subsubsection{水质与环境观察}

调研中发现鱼塘存在明显的水质问题：
\begin{itemize}[leftmargin=2em]
    \item 鱼塘水面出现富营养化现象，藻类密集覆盖水面；
    \item 增氧机为唯一标配设备，智能化程度较低。
\end{itemize}

\begin{figure}[H]
\centering
\fbox{\parbox{0.8\textwidth}{\centering\vspace{1.5cm}【图 北定福村调研现场】\vspace{1.5cm}}}
\caption{北定福村调研现场}
\end{figure}

% ---------- 3.2 东双营村 ----------
\subsection{东双营村渔业调研}

东双营村主要产业已转向宠物文化产业，但仍保留约250亩鱼塘（7块水域），由村干部带头建设于山区。

\subsubsection{养殖模式与规模}

\begin{itemize}[leftmargin=2em]
    \item 鱼塘面积约250亩，共7块水域；
    \item 由村干部带头建设，位于山区；
    \item 依赖地下水，几乎不换水；
    \item 受雨水影响大，没有形成系统的养殖体系。
\end{itemize}

\subsubsection{智能化探索}

养殖户老赵为提升养殖品质，添置了智能设备，总投入超过200万元。但调研发现"纯利不高"，高投入并未带来相应的高回报。7--8月份需勤溜达、高要求管理，养殖周期约23天。

\subsubsection{销售渠道}

销售完全依赖鱼贩子收购，无网络销售渠道。鱼价约7--9元/斤，价格不稳定。没有通过网上直播售卖鱼。

\subsubsection{养殖管理}

\begin{itemize}[leftmargin=2em]
    \item 有增加高价种类的鱼；
    \item 鱼病以烂鳃为主，一般是细菌；
    \item 养鱼的人变少了，年轻人不愿意干这个了；
    \item 缺钱，缺宣传；
    \item 交通特别方便，有望联动来做旅游业；
    \item 35亩建设用地（2020）不了了之；
    \item 鱼病了要药品，要上实验室。
\end{itemize}

\begin{figure}[H]
\centering
\fbox{\parbox{0.8\textwidth}{\centering\vspace{1.5cm}【图 东双营村鱼塘】\vspace{1.5cm}}}
\caption{东双营村鱼塘}
\end{figure}

% ---------- 3.3 天井村 ----------
\subsection{天井村渔业调研}

天井村是调研村庄中生态模式最为先进的村庄，其水质管理经验具有重要的推广价值。

\subsubsection{养殖模式与规模}

\begin{itemize}[leftmargin=2em]
    \item 全村有五六户养殖户，共百十亩鱼塘，每个30--40亩；
    \item 分散经营，无闲置塘；
    \item 个人承包模式。
\end{itemize}

\subsubsection{水质管理亮点}

天井村最大的亮点在于水质管理，形成了独特的"鱼---藕"生态闭环：
\begin{itemize}[leftmargin=2em]
    \item \textbf{光合细菌调水}：使用光合细菌调节水质，增加水体含氧量；
    \item \textbf{尾水循环利用}：尾水排入藕池循环利用，形成生态闭环；
    \item \textbf{藕塘净化}：村里有两个专门的藕塘（因渗水比较严重），在净化功能之外额外产出莲藕；
    \item 夏季养鱼氧气多，没开增氧机。
\end{itemize}

\subsubsection{销售渠道创新}

天井村是调研村庄中唯一有微信本地销售渠道的村庄。部分养殖户通过"老板收货+微信导鱼"的双渠道模式获得更高收益。另有企业对接，收益更高一点。

\subsubsection{其他观察}

\begin{itemize}[leftmargin=2em]
    \item 草鱼、鲤鱼亩产约10000斤，产量较高；
    \item 鱼苗来源多样，价格各异，孵化技术依赖经验，成活率看眼力；
    \item 一年收成不确定，约7--8元一斤；
    \item 一个品种一个养殖方式，困难；
    \item 有环保要求，不许排污地下水；
    \item 政府没有补贴，自负盈亏，税收不收税；
    \item 鱼病有腮病、蛆虫（鱼表面不光滑）等问题，需请鱼大夫或去农科院；
    \item 突然性断电需自备发电机；
    \item 有手机联网等高科技手段辅助管理；
    \item 观赏鱼养殖，翠绿色最佳；
    \item 冬天增加水位养鱼，周期比较短；
    \item 下雨影响大，鸟常常叼鱼走；一般不换水（返水需换，常年下雨）。
\end{itemize}

\begin{figure}[H]
\centering
\fbox{\parbox{0.8\textwidth}{\centering\vspace{1.5cm}【图 天井村鱼塘渔船作业】\vspace{1.5cm}}}
\caption{天井村鱼塘渔船作业}
\end{figure}

% ---------- 3.4 马昌营村 ----------
\subsection{马昌营村渔业调研}

马昌营村是镇中心村，鱼塘数量最多，是调研中唯一"全品种覆盖"的村庄。

\subsubsection{养殖模式与规模}

\begin{itemize}[leftmargin=2em]
    \item 鱼塘数量最多，共十六七个；
    \item 原有七户散户养殖，部分已不再经营，鱼塘正逐步集中；
    \item 同时养殖食用鱼和观赏鱼，为调研中唯一"全品种覆盖"的村庄；
    \item 租赁土地经营。
\end{itemize}

\subsubsection{经营困境}

养殖户老侯详细描述了经营困境：
\begin{itemize}[leftmargin=2em]
    \item \textbf{销售期短}：一年只能养一茬，1亩/2000多斤，冬天水结冰卖不掉；
    \item \textbf{设备老旧}：增氧机已使用六七年，早年有政府补贴，如今纯自费维护；
    \item \textbf{智能设备缺失}：自动投喂机单价3000至4000元，"太贵了"没有购置；
    \item \textbf{投喂方式原始}：目前为手动撒药、手动投喂。
\end{itemize}

\subsubsection{销售痛点}

"专业贩子来收，贩子往网上挂，大小都收。"贩子以7至9元/斤收购后在网上以15至25元销售。当被问及是否愿意网上卖鱼时，老侯回答："那敢情好。"

\subsubsection{水质与病害}

\begin{itemize}[leftmargin=2em]
    \item 水质不好，一下雨需要调水质；
    \item 试纸测水质，养殖户自行判断；
    \item 鱼苗自己孵化+购买，成活率约90\%；
    \item 鱼病以鱼腮病、吃多为主；
    \item 药品由公司供应，大病少、小病多；
    \item 尾水导着使，一般不发水；
    \item 成本高收入少。
\end{itemize}

\begin{figure}[H]
\centering
\fbox{\parbox{0.8\textwidth}{\centering\vspace{1.5cm}【图 马昌营村养殖场全景】\vspace{1.5cm}}}
\caption{马昌营村养殖场全景（增氧机、防鸟网、作业船）}
\end{figure}

% ---------- 3.5 王各庄村 ----------
\subsection{王各庄村渔业调研}

王各庄村是第二大村，历史悠久，面积2.49平方公里，约500户，有特殊的盘窑技术，一直在申请非遗（3位技术人）。鱼塘总面积约400亩，分散经营，全部养殖食用鱼，渔业为非主营产业。

\subsubsection{养殖模式与规模}

\begin{itemize}[leftmargin=2em]
    \item 鱼塘总面积约400亩，分散经营，非主营；
    \item 全部养殖食用鱼，无观赏鱼；
    \item 特点：隔开一个个小池塘；
    \item 部分闲置鱼塘已改种莲藕（5--6个），尝试生态转型；
    \item 不换水，不让导水；尾水禁止排，排地里。
\end{itemize}

\subsubsection{第一位受访养殖户}

\begin{itemize}[leftmargin=2em]
    \item \textbf{经营困境}："最近亏钱，市场不行""销量从2万变5000""养鱼不挣钱，导鱼的挣钱"；
    \item \textbf{鱼苗问题}：鱼苗坐飞机来，成活率低，60\%不到；
    \item \textbf{智能设备}：使用手机"塘管家"APP管理水质，年费100多元，但"岁数大的不会用"，准确度高；
    \item \textbf{用药控制}：用药控制严格，鱼大夫（只卖药）自己请，一般自己搞；
    \item \textbf{鱼病}：裂头病；
    \item \textbf{饲料成本}：约1万元/吨；
    \item \textbf{年纯收入}：仅5至6万元；
    \item \textbf{其他}：政府没有补贴，扩大不了（没地方了），只有考不上的大学生来养鱼；没有停电停水；一年只有几万块钱；成熟周期1年；希望不要有鱼病；销售不出去；平谷水质好；网上直播卖鱼不会操作，让下一辈；有自动投食的机器。
\end{itemize}

\subsubsection{第二位受访养殖户}

\begin{itemize}[leftmargin=2em]
    \item 20亩，2个池，食用鱼（草鱼）；
    \item 一年一茬，6万斤产量；
    \item 20多年养鱼经验；
    \item 鱼苗自己培养的，外地买的无籽，成活率草鱼最高70\%；
    \item 夏季鱼病频繁，烂鳃为主；
    \item 只有增氧机，无其他智能设备；有发电机，几分停电停水；
    \item 有人帮忙测水质；用多的水果来做鱼饵；
    \item 全是鱼贩子销售；
    \item 尾水往河道排，如果有玉米之类可循环处理（暂未养）；
    \item 政府无补贴和帮助；
    \item 养鱼不挣钱，导鱼的挣钱；能挣5--6万一年。
\end{itemize}

\begin{figure}[H]
\centering
\fbox{\parbox{0.8\textwidth}{\centering\vspace{1.5cm}【图 王各庄村鱼塘】\vspace{1.5cm}}}
\caption{王各庄村鱼塘}
\end{figure}

% ---------- 3.6 后芮营村（芮朝利） ----------
\subsection{后芮营村（芮朝利）渔业调研}

后芮营村是本次调研中智能化程度最高、产业模式最先进的村庄。养殖户芮朝利采用"鱼菜共生"生态模式，实现了全年养殖和较高经济效益。

\subsubsection{养殖模式与规模}

\begin{itemize}[leftmargin=2em]
    \item 总共50亩，水面38.8亩；
    \item 鱼菜共生模式（水体过剩），搭配水稻、荷花环保种植；
    \item 主业养鱼，一家人经营（未与村民合伙，因多户人话语权分散）；
    \item 有合作社"源来"。
\end{itemize}

\subsubsection{智能化设备配置}

该村是调研中智能化程度最高的：
\begin{itemize}[leftmargin=2em]
    \item \textbf{手机管理增氧机}：使用"塘管家"APP，定时控制，人不用巡塘；
    \item \textbf{投饲机}：自动投喂，减少鱼病；
    \item \textbf{拉网机}：辅助捕捞；
    \item \textbf{催水设备}：催水鱼养殖；
    \item \textbf{大棚养殖}：政府有补贴，100--200万；
    \item \textbf{每个鱼塘4个发电机}：保障电力供应。
\end{itemize}

\subsubsection{品种与产量}

\begin{itemize}[leftmargin=2em]
    \item 主养鲈鱼（为主），18--19元/斤批发价；
    \item 五彩色的鱼（锦鲤）、黑色的鱼不易被吃，所以有特殊养殖策略；
    \item 产量：100吨；
    \item 利润：20--30万（受大鱼病影响大）；
    \item 苗期自己培养鱼苗，增加鱼收，但具有市场滞后性。
\end{itemize}

\subsubsection{全年养殖模式}

\begin{itemize}[leftmargin=2em]
    \item 一般养鱼时间是4--10月；
    \item 增加养殖周期，实现全年养殖；
    \item 冬天不能使用大棚（可能指部分设备）；
    \item 每年抽一部分水；
    \item 预备进行尾水治理实验过程。
\end{itemize}

\subsubsection{销售与成本}

\begin{itemize}[leftmargin=2em]
    \item 批发价为主，零售麻烦；
    \item 价格看市场，有时14元，价格过低；
    \item 饲料占大头；
    \item 世面上的拉网机、撒料船和风投等高端设备不太稳定。
\end{itemize}

\subsubsection{鱼病防治与技术输出}

\begin{itemize}[leftmargin=2em]
    \item 鱼病主要是鱼鳃；
    \item 帮助别人治鱼病，半义诊，钱一般收药钱；
    \item 有讲课、定期看水的地方，自己会看一些；
    \item 带动作用：义诊、参观讲课，愿意给技术性输出；
    \item 有考虑AI、机器人之类的引入（我们的构思：报错后，机器人去）。
\end{itemize}

\subsubsection{发展瓶颈}

\begin{itemize}[leftmargin=2em]
    \item 不适合散户（大棚模式需要较大投入）；
    \item 需要招商，愿意继续完善信息化展示网站规划，愿意给更多信息；
    \item 目前不是特别需要招商，但持开放态度。
\end{itemize}

\begin{figure}[H]
\centering
\fbox{\parbox{0.8\textwidth}{\centering\vspace{1.5cm}【图 后芮营村智能化养殖大棚】\vspace{1.5cm}}}
\caption{后芮营村智能化养殖大棚}
\end{figure}

% ---------- 3.7 西海子村 ----------
\subsection{西海子村渔业调研}

西海子村以观赏鱼养殖为特色，是调研中品种最丰富、养殖规模较大的村庄之一。

\subsubsection{养殖模式与规模}

\begin{itemize}[leftmargin=2em]
    \item 8个池塘，55亩，养殖面积30多亩；
    \item 观赏鱼（合伙养殖），淡水鱼赔钱；
    \item 独立养殖，有合作社"源来"；
    \item 孙老师负责，使用地下水。
\end{itemize}

\subsubsection{品种构成}

\begin{itemize}[leftmargin=2em]
    \item 黄金鲤：100多万条；
    \item 金鱼：1万多条；
    \item 锦鲤：产量不太高；
    \item 等等，品种很多。
\end{itemize}

\subsubsection{产量与销售}

\begin{itemize}[leftmargin=2em]
    \item 产量：5--6万斤；
    \item 养殖业价格不稳定；
    \item 运输繁琐，网络不搞了，一般给市场；
    \item 秋后至来年开春，观赏鱼随时都有人要；
    \item 现在采用智能and半智能养殖。
\end{itemize}

\subsubsection{成本与设备}

\begin{itemize}[leftmargin=2em]
    \item 饲料30--50万，太贵了（赔钱原因），3元/斤，还有时段影响；
    \item 设备一般，缺陷较大；
    \item 一天700多度电，电费也贵；
    \item 水质试剂测，鱼大大（敏感度高）。
\end{itemize}

\subsubsection{鱼病与风险}

\begin{itemize}[leftmargin=2em]
    \item 鱼病主要是鱼鳃；
    \item 鱼病注意防治，调水；
    \item 有讲课，定期看水的地方，自己会看一些；
    \item 村里没有鱼医；
    \item 有可能形成局部龙卷风，会造成经济损伤；
    \item 养殖一般没保险，风险太高且不可估量；保费高，不像种植的一般交费低且保额高。
\end{itemize}

\subsubsection{尾水与文旅设想}

\begin{itemize}[leftmargin=2em]
    \item 尾水排到坑里，专门搞了一个坑（里面有鱼）；
    \item 有文旅想法，政策不好搞，建设用地和文旅用地不同；
    \item 没有配套的，需要板块合并；
    \item 想要农场一日游，没有人牵头。
\end{itemize}

\begin{figure}[H]
\centering
\fbox{\parbox{0.8\textwidth}{\centering\vspace{1.5cm}【图 西海子村观赏鱼养殖池塘】\vspace{1.5cm}}}
\caption{西海子村观赏鱼养殖池塘}
\end{figure}

% ==================== 四、渔村调研数据分析 ====================
\section{渔村调研数据分析}

\subsection{七村渔业调研对比}

通过对七个村庄的交叉比对，整理调研数据如下：

\begin{table}[H]
\centering
\caption{七村渔业调研对比表}
\small
\begin{tabular}{>{\centering\arraybackslash}p{1.6cm}|>{\centering\arraybackslash}p{1.5cm}|>{\centering\arraybackslash}p{1.3cm}|>{\centering\arraybackslash}p{1.3cm}|>{\centering\arraybackslash}p{1.3cm}|>{\centering\arraybackslash}p{1.3cm}|>{\centering\arraybackslash}p{1.3cm}|>{\centering\arraybackslash}p{1.3cm}}
\toprule
\textbf{项目} & \textbf{北定福} & \textbf{东双营} & \textbf{天井村} & \textbf{马昌营} & \textbf{王各庄} & \textbf{后芮营} & \textbf{西海子} \\
\midrule
养殖模式 & 集中连片 & 山区干部带头 & 老板+散户 & 16--17塘 & 分散非主营 & 鱼菜共生 & 观赏鱼合伙 \\
\midrule
规模 & 中小型 & 250亩/7块 & 百十亩 & 最大 & 400亩 & 50亩/38.8水面 & 55亩/8塘 \\
\midrule
主要品种 & 鲤鱼草鱼 & 综合+高价 & 综合鱼种 & 食用+观赏 & 食用鱼 & 鲈鱼为主 & 黄金鲤金鱼 \\
\midrule
亩产 & 约2000斤 & 约2000斤 & 约10000斤 & 约2000斤 & 约2000斤 & 100吨总产 & 5--6万斤 \\
\midrule
销售渠道 & 鱼贩子 & 鱼贩子 & 老板+微信 & 鱼贩子 & 鱼贩子 & 批发为主 & 市场 \\
\midrule
智能设备 & 增氧机 & 智能设备(200万) & 增氧机 & 增氧机 & 塘管家APP & 手机+投饲+拉网 & 试剂测水 \\
\midrule
水质管理 & 常规 & 几乎不换水 & 光合细菌+藕池 & 试纸检测 & 有人测水质 & 手机定时管理 & 试剂测+鱼大大 \\
\midrule
年利润 & 微利 & 纯利不高 & 不稳定 & 成本高收入少 & 5--6万 & 20--30万 & 亏钱 \\
\bottomrule
\end{tabular}
\end{table}

\subsection{核心发现}

\subsubsection{发现一：鱼价十年停滞，利润被中间商截取}

各村庄鱼价每斤7至9元（后芮营鲈鱼18--19元为特例），与十年前基本持平。贩子以7至9元收购，网上卖15至25元，差价全部被中间环节截取。养殖户辛苦一年，利润却被中间商大幅吞噬。

\subsubsection{发现二：从业人员全面老龄化，后继无人}

受访养殖户全部50岁以上。"年轻人不愿意干""只有考不上的大学生来养鱼"——七个村庄反复出现这一困境。渔业面临严重的人才断层。

\subsubsection{发现三：设备更新滞后，智能化程度两极分化}

\begin{itemize}[leftmargin=2em]
    \item \textbf{低端}：北定福、马昌营、王各庄仅增氧机，自动投喂机3000至4000元无人购置，水质检测靠试纸，鱼病靠肉眼观察；
    \item \textbf{高端}：后芮营村实现手机定时管理增氧机、自动投饲机、拉网机、大棚养殖，政府补贴100--200万；
    \item \textbf{中间}：东双营投入200万智能设备但"纯利不高"，王各庄使用"塘管家"APP但老人不会用。
\end{itemize}

\subsubsection{发现四：天井村生态模式具有推广价值}

光合细菌调水+藕池循环利用的"鱼---藕"闭环，成本低、易复制，是唯一具备系统水质管理方案的传统村庄。该模式可在全镇范围内推广。

\subsubsection{发现五：后芮营村智能化模式代表未来方向}

后芮营村的鱼菜共生、大棚养殖、手机远程管理、自动投饲等模式，实现了全年养殖和20--30万年利润，是调研中经济效益最好的案例。但该模式需要较大投入（政府补贴100--200万），不适合散户复制。

\subsubsection{发现六：养殖户愿意改变但缺乏路径指导}

老赵投入智能设备、老侯推动散户整合、天井村构建生态循环、芮朝利探索AI机器人——他们一直在寻找出路，但尝试是零散的、自发的。老侯那句"那敢情好"，是最真实的需求表达。

\subsubsection{发现七：宠物产业链为渔业转型提供参照}

东双营宠物产业园已形成"生产---仓储---直播---物流"完整闭环，单场直播最高300万元。同处一地，渔业完全可以复制"产地直发+直播电商"模式。

\subsubsection{发现八：观赏鱼市场存在结构性机会}

西海子村黄金鲤100多万条、金鱼1万多条的规模，以及"秋后至来年开春随时有人要"的特性，说明观赏鱼市场有稳定需求，但受饲料成本（30--50万）和电费（一天700多度）拖累，整体亏损。

% ==================== 五、特色产业与文化活动调研 ====================
\section{特色产业与文化活动调研}

除渔村主线调研外，实践队还参观了镇域内的特色产业项目和文化活动，了解马昌营镇的多元产业生态。

\subsection{宠物文化产业园参观}

实践队参观了位于东双营的北京（平谷）宠物文化产业园——PETIDEA宠想里。园区汇聚了宠物食品用品生产企业、宠物主题商业空间和大型仓储物流中心，形成了"生产---仓储---直播带货---物流配送"的完整产业链闭环。

据了解，园区内商家通过直播带货，单场销售额可达50万元，高峰期甚至突破300万元。这一"邻居产业"的成功经验为渔业转型提供了重要参照——同处一地，宠物产业已实现数字化全链路运营，而渔业仍停留在最原始的"养殖---贩子---市场"环节。

\begin{figure}[H]
\centering
\fbox{\parbox{0.8\textwidth}{\centering\vspace{1.5cm}【图 PETIDEA宠想里宠物文化产业园】\vspace{1.5cm}}}
\caption{PETIDEA宠想里宠物文化产业园}
\end{figure}

\subsection{"又见咖啡"乡村文旅调研}

实践队参观了马昌营镇"又见咖啡"乡村文旅项目，了解了乡村休闲产业与咖啡文化的融合探索。该项目将闲置农房改造为特色咖啡空间，为乡村带来了新的消费场景和年轻客流。这种"内容驱动"的文旅模式为后续提出休闲渔业建议提供了参考。

\begin{figure}[H]
\centering
\fbox{\parbox{0.8\textwidth}{\centering\vspace{1.5cm}【图 "又见咖啡"乡村文旅项目】\vspace{1.5cm}}}
\caption{"又见咖啡"乡村文旅项目}
\end{figure}

\subsection{非遗集中活动}

实践队参加马昌营镇政府组织的非遗集中活动，深入了解了马昌营镇景泰蓝杯垫制作等非物质文化遗产的传承现状。团队成员与非遗传承人面对面交流，亲手体验了景泰蓝工艺流程。

\begin{figure}[H]
\centering
\fbox{\parbox{0.8\textwidth}{\centering\vspace{1.5cm}【图 景泰蓝杯垫制作非遗活动现场】\vspace{1.5cm}}}
\caption{景泰蓝杯垫制作非遗活动现场}
\end{figure}

\subsection{物流中心参观}

实践队还参观了镇域内的宠物用品仓储物流中心，了解了"直播+仓储+物流"的电商基础设施。

\subsection{"首爱"集中活动}

实践队在天井村参加"首爱"集中活动，与首都爱护动物协会的志愿者们共同开展了流浪动物保护宣传和公益领养活动。实践队员协助志愿者进行场地布置、信息登记和宣传讲解工作，体验了"产业+公益"双轮驱动的乡村服务模式。

% ==================== 六、实践成果汇总 ====================
\section{实践成果汇总}

截至2026年7月19日，"数智渔韵"实践队已取得以下阶段性成果：

\begin{table}[H]
\centering
\caption{实践成果汇总表}
\begin{tabular}{clp{8cm}}
\toprule
\textbf{序号} & \textbf{成果类型} & \textbf{成果说明} \\
\midrule
1 & 大报告 & 本成果报告文档（覆盖7个村庄调研） \\
2 & 每日工作日志 & 每日撰写实践日志，记录当日调研内容与心得 \\
3 & 视频 & 《我们都是追梦人》短视频、每周亮点视频 \\
4 & 招商网站 & 马昌营镇对外招商网站（已上线） \\
5 & 公众号 & "数智渔韵"微信公众号，已发布多篇推文 \\
6 & 调研素材 & 照片200余张，访谈视频60余条 \\
\bottomrule
\end{tabular}
\end{table}

% ==================== 七、意见建议 ====================
\section{意见建议}

基于六天的实地调研，实践队从智能科学与技术专业视角出发，提出以下建议：

\subsection{建立"直播+仓储"渔业电商体系}

参照宠物产业园"单场直播50万至300万"的成功经验，由专业团队负责品牌设计、直播运营和物流配送，养殖户专注把鱼养好。可参考"平谷大桃"的品牌化经验，打造"平谷鱼"区域公用品牌。

\subsection{推广天井村"鱼---藕"生态循环模式}

将光合细菌调水+藕池循环利用的技术方案标准化，通过区农业部门组织培训、制作短视频教程、在示范鱼塘安装简易水质监测设备等方式推广。

\subsection{开展智能投喂与AI鱼病诊断试点}

自动投喂机日均成本仅2至3.5元，建议由高校或政府提供设备供一个鱼塘试用。同时可引入基于图像识别的AI鱼病诊断、智能水质监测等技术在示范点先行先试。后芮营村已有手机管理增氧机的成功经验，可作为试点基础。

\subsection{探索休闲渔业与文旅融合}

结合平谷桃花节、大桃采摘等成熟IP和"又见咖啡"等已有文旅项目，打造"夏赏荷花秋品鱼"的全季旅游产品。西海子村的观赏鱼养殖可发展为观光项目。

\subsection{建立"高校+政府+养殖户"长效合作机制}

依托高校的智能技术能力，为渔业转型提供持续的技术支持。实践队开发的招商网站和公众号可以作为校地合作的数字化纽带。后芮营村愿意接受技术性输出和AI机器人引入，可作为首个合作试点。

\subsection{推动观赏鱼产业差异化发展}

西海子村黄金鲤100多万条的规模具有市场优势，建议：
\begin{itemize}[leftmargin=2em]
    \item 降低饲料成本（目前30--50万/年，3元/斤）；
    \item 解决高电费问题（一天700多度）；
    \item 发展"观赏鱼+文旅"模式，实现农场一日游；
    \item 引入养殖保险，降低龙卷风等自然灾害风险。
\end{itemize}

\end{document}
